\subsection{弗雷德里克·里斯理论}

设 E 是分离局部凸空间，h 是 E 的一个紧自同态。设 \(\lambda \in K\) 。自同态 \(\lambda h\) 是紧的，故 \(1_{E} - \lambda h\) 是 E 的里斯自同态（III, p. 75 的命题 2 的推论 2）；用 \(N_{\lambda}\) 与 \(I_{\lambda}\) 表示它的幂零空间与余幂零空间。由 III, p. 47 的命题 6，我们有下列性质：

\begin{enumerate}
\def\labelenumi{(\roman{enumi})}
\item
  E 的向量子空间 \(I_{\lambda}\) 与 \(N_{\lambda}\) 是闭的且在 h 下稳定；
\item
  局部凸空间 E 是 \(I_{\lambda}\) 与 \(N_{\lambda}\) 的拓扑直和；
\item
  映射 \(1_{\mathrm{E}} - \lambda h\) 通过限制诱导 \(\mathrm{I}_{\lambda}\) 的一个自同构；
\item
  向量空间 \(N_{\lambda}\) 有限维，且存在一个整数 \(n_{\lambda} \geqslant 0\) ，使得 \((1_{\mathrm{E}} - \lambda h)^{n_{\lambda}}(x) = 0\) （对一切 \(x \in N_{\lambda}\) ）。
\end{enumerate}

这些性质以唯一的方式确定空间 \(I_{\lambda}\) 与 \(N_{\lambda}\) （A, VIII, p. 26, 注记 2）。

引理 3. --- 设 X 是一个拓扑具有可数基的拓扑空间。X 的每个离散子集都是可数的。

设 \(\mathcal{B}\) 是 X 的拓扑的一个可数基，D 是 X 的一个离散子集。对 D 的每个元素 \(x\) ，存在元素 \(\mathrm{B}_x\) （属于 \(\mathcal{B}\) ），使得 \(\mathrm{B}_x \cap \mathrm{D} = \{x\}\) 。从 D 到 \(\mathcal{B}\) 的映射由 \(x \mapsto \mathrm{B}_x\) 定义，它是单射，由此引理得证。

定理 4. --- 设 E 是分离局部凸空间，h 是 E 的一个紧自同态。集 \(\Sigma\) 由那些 \(\lambda \in K\) （使 \(1_{\mathrm{E}} - \lambda h\) 不是 E 的自同构的）组成，它是域 K 的一个可数、闭且离散的子集。

集 \(\Sigma\) 也是由那些 \(\lambda\in\mathrm{K}\) （使 \(1_{\mathrm{E}}-\lambda h\) 不是单射的）所成的集，以及由那些 \(\lambda\in\mathrm{K}\) （使 \(1_{\mathrm{E}}-\lambda h\) 不是满射的）所成的集。

存在一个巴拿赫空间 G、一个连续线性映射 \(p: E \rightarrow G\) 和一个紧线性映射 \(q: G \rightarrow E\)，使得 \(h = q \circ p\) （III, p. 5, 命题 5）。G 的自同态 \(h' = p \circ q\) 是紧的，且 \(1_{E} - \lambda h\) 是 E 的自同构当且仅当 \(1_{G} - \lambda h'\) 是 G 的自同构（III, p. 49 的命题 10, e)）。因此我们只需在 E 是巴拿赫空间的情形证明该定理，以下即作此假设。

设 \(\lambda \in K\) 。由于里斯自同态的指标为零，下列说法等价：映射 \(1_{\mathrm{E}} - \lambda h\) 是 E 的自同构，或它是单射的，或它是满射的。

集 \(\Sigma\) 在 K 中是闭的，因为 E 的自同构所成的集在 \(\mathcal{L}(\mathrm{E})\) 中是开的。设 \(\lambda\) 是 \(\Sigma\) 的一个元素。我们有 \(\lambda \neq 0\) 。用 N 表示 h 的幂零空间，I 表示其余幂零空间，并用 \(h_{I}\) 与 \(h_{N}\) 表示 h 在这些子空间上诱导的 I 与 N 的自同态。则 \(1_{I} - \lambda h_{I}\) 是 I 的一个自同构，且存在 \(\lambda\) 在 K 中的一个邻域 U，使得 \(1_{I} - \mu h_{I}\) 是 I 的自同构（对一切 \(\mu \in U\) ）。N 的自同态 \(1_{N} - \lambda h_{N}\) 是幂零的；对一切 \(\mu \neq \lambda\) ，我们有

\[
1 _ {\mathrm{N}} - \mu h _ {\mathrm{N}} = \frac {\lambda - \mu}{\lambda} \Big (1 _ {\mathrm{N}} + \frac {\mu}{\lambda - \mu} (1 _ {\mathrm{N}} - \lambda h _ {\mathrm{N}}) \Big),
\]

故 \(1_{N} - \mu h_{N}\) 是 N 的一个自同构。于是集 \(U \cap \Sigma\) 缩减为单个元素 \(\lambda\) ，从而集 \(\Sigma\) 是离散的。由引理 3 它是可数的。

